\documentclass[aps,preprint]{revtex4}%
\usepackage{amssymb}
\usepackage{amsfonts}%
\usepackage{amsmath}%
\setcounter{MaxMatrixCols}{30}%
\usepackage{graphicx}
%TCIDATA{OutputFilter=latex2.dll}
%TCIDATA{Version=5.50.0.2960}
%TCIDATA{CSTFile=revtex4.cst}
%TCIDATA{Created=Wednesday, April 27, 2016 09:13:10}
%TCIDATA{LastRevised=Wednesday, April 27, 2016 09:36:02}
%TCIDATA{<META NAME="GraphicsSave" CONTENT="32">}
%TCIDATA{<META NAME="SaveForMode" CONTENT="1">}
%TCIDATA{BibliographyScheme=Manual}
%TCIDATA{<META NAME="DocumentShell" CONTENT="Articles\SW\REVTeX 4">}
%BeginMSIPreambleData
\providecommand{\U}[1]{\protect\rule{.1in}{.1in}}
%EndMSIPreambleData
\newtheorem{theorem}{Theorem}
\newtheorem{acknowledgement}[theorem]{Acknowledgement}
\newtheorem{algorithm}[theorem]{Algorithm}
\newtheorem{axiom}[theorem]{Axiom}
\newtheorem{claim}[theorem]{Claim}
\newtheorem{conclusion}[theorem]{Conclusion}
\newtheorem{condition}[theorem]{Condition}
\newtheorem{conjecture}[theorem]{Conjecture}
\newtheorem{corollary}[theorem]{Corollary}
\newtheorem{criterion}[theorem]{Criterion}
\newtheorem{definition}[theorem]{Definition}
\newtheorem{example}[theorem]{Example}
\newtheorem{exercise}[theorem]{Exercise}
\newtheorem{lemma}[theorem]{Lemma}
\newtheorem{notation}[theorem]{Notation}
\newtheorem{problem}[theorem]{Problem}
\newtheorem{proposition}[theorem]{Proposition}
\newtheorem{remark}[theorem]{Remark}
\newtheorem{solution}[theorem]{Solution}
\newtheorem{summary}[theorem]{Summary}
\newenvironment{proof}[1][Proof]{\noindent\textbf{#1.} }{\ \rule{0.5em}{0.5em}}
\begin{document}
\preprint{HEP/123-qed}
\title[Short title for running header]{REV\TeX4}
\author{First Author}
\affiliation{My Institution}
\author{Second Author}
\affiliation{My Institution}
\author{Third Author}
\affiliation{Other Institution}
\keywords{one two three}
\pacs{PACS number}

\begin{abstract}
Shell document for REV\TeX{} 4.

\end{abstract}
\volumeyear{year}
\volumenumber{number}
\issuenumber{number}
\eid{identifier}
\date[Date text]{date}
\received[Received text]{date}

\revised[Revised text]{date}

\accepted[Accepted text]{date}

\published[Published text]{date}

\startpage{101}
\endpage{102}
\maketitle
\tableofcontents


\subsubsection{Characteristic Functions (Generatoring Functions)}

\paragraph{Majoranna Form of Geminal Operators}

Note that for \emph{real} geminal operators $g$ i.e. elements of
$\mathcal{F}_{2}$ one has
\begin{align}
g^{t} &  =-g;g^{\dagger}=-g\\
\left(  e^{g}\right)  ^{t} &  =e^{g^{t}}=e^{-g};\left(  e^{g}\right)
^{\dagger}=e^{g^{\dagger}}=e^{-g},
\end{align}
but note that
\begin{equation}
\left(  ig\right)  ^{\dagger}=ig
\end{equation}
i.e. $ig$ is self adjoint.

However in general%
\begin{align}
\left[  g_{1},g_{2}\right]    & \neq0;\left[  g_{3},\left[  g_{1}%
,g_{2}\right]  \right]  \neq0\\
e^{g_{1}+g_{2}}  & \neq e^{g_{1}}e^{g_{2}}.
\end{align}


$\lceil$
\begin{align}
x_{1}x_{2}x_{1}x_{3}  & =-x_{2}x_{3}\\
x_{1}x_{3}x_{1}x_{2}  & =-x_{3}x_{2}=x_{2}x_{3}\\
\left[  x_{1}x_{2},x_{1}x_{3}\right]    & =-2x_{2}x_{3}\neq0.
\end{align}
$\rfloor$ 

One can define two characteristic function, $\Phi_{rD},$ $\Phi_{iD}$
associated with $D$ given by
\begin{align}
\Phi_{rD} &  =\mathit{\operatorname*{tr}}\left(  \exp\left\{  g\right\}
D\right)  \in\mathbb{R}\\
\Phi_{iD} &  =\mathit{\operatorname*{tr}}\left(  \exp\left\{  ig\right\}
D\right)  \in\mathbb{C}.
\end{align}
We can always express the geminal in canonical form as
\begin{equation}
g=%
%TCIMACRO{\tsum \limits_{1\leq j\leq r}}%
%BeginExpansion
{\textstyle\sum\limits_{1\leq j\leq r}}
%EndExpansion
\mathit{g}_{yjj+r}y_{j}y_{j+r};\mathit{g}_{y1jj+r}\in\mathbb{R}.
\end{equation}


The expansion of the real exponential gives
\begin{equation}
\exp\left\{  g\right\}  =I+g+\frac{1}{2!}g^{2}+\frac{1}{3!}g^{3}+\cdots
=\exp\left\{
%TCIMACRO{\tsum \limits_{1\leq j\leq r}}%
%BeginExpansion
{\textstyle\sum\limits_{1\leq j\leq r}}
%EndExpansion
\mathit{g}_{yjj+r}y_{j}y_{j+r}\right\}  ,
\end{equation}
which can be developed to give
\begin{align}
&  \exp\left\{
%TCIMACRO{\tsum \limits_{1\leq j\leq r}}%
%BeginExpansion
{\textstyle\sum\limits_{1\leq j\leq r}}
%EndExpansion
\mathit{g}_{yjj+r}y_{j}y_{j+r}\right\}  =%
%TCIMACRO{\tprod \limits_{1\leq j\leq r}}%
%BeginExpansion
{\textstyle\prod\limits_{1\leq j\leq r}}
%EndExpansion
\exp\left\{  \mathit{g}_{yjj+r}y_{j}y_{j+r}\right\}  \\
&  =%
%TCIMACRO{\tprod \limits_{1\leq j\leq r}}%
%BeginExpansion
{\textstyle\prod\limits_{1\leq j\leq r}}
%EndExpansion
\left\{  I_{\mathcal{F}}+\mathit{g}_{yjj+r}y_{j}y_{j+r}+\frac{1}{2}%
\mathit{g}_{j_{1}j_{2}}^{2}\left(  y_{j}y_{j+r}\right)  ^{2}+\frac{1}%
{3!}\mathit{g}_{j_{1}j_{2}}^{3}\left(  y_{j}y_{j+r}\right)  ^{3}%
+\cdots\right\}  \\
&  =%
%TCIMACRO{\tprod \limits_{1\leq j\leq r}}%
%BeginExpansion
{\textstyle\prod\limits_{1\leq j\leq r}}
%EndExpansion
\left\{  I_{\mathcal{F}}+\mathit{g}_{yjj+r}y_{j}y_{j+r}-\frac{1}{2}%
\mathit{g}_{yjj+r}^{2}I_{\mathcal{F}}-\frac{1}{3!}\mathit{g}_{yjj+r}^{3}%
y_{j}y_{j+r}+\frac{1}{4!}\mathit{g}_{yjj+r}^{4}I_{\mathcal{F}}+\frac{1}%
{5!}\mathit{g}_{yjj+r}^{5}y_{j}y_{j+r}+\cdots\right\}  \\
&  =%
%TCIMACRO{\tprod \limits_{1\leq j\leq r}}%
%BeginExpansion
{\textstyle\prod\limits_{1\leq j\leq r}}
%EndExpansion
\left\{  I_{\mathcal{F}}-\frac{1}{2}\mathit{g}_{yjj+r}^{2}I_{\mathcal{F}%
}+\frac{1}{4!}\mathit{g}_{yjj+r}^{4}I_{\mathcal{F}}+\mathit{g}_{yjj+r}%
y_{j}y_{j+r}-\frac{1}{3!}\mathit{g}_{yjj+r}^{3}y_{j}y_{j+r}+\frac{1}%
{5!}\mathit{g}_{yjj+r}^{5}y_{j}y_{j+r}-\cdots\right\}  \\
&  =%
%TCIMACRO{\tprod \limits_{1\leq j\leq r}}%
%BeginExpansion
{\textstyle\prod\limits_{1\leq j\leq r}}
%EndExpansion
\left\{  \left(  1-\frac{1}{2}\mathit{g}_{yjj+r}^{2}+\frac{1}{4!}%
\mathit{g}_{yjj+r}^{4}-+\cdots\right)  I_{\mathcal{F}}+\left(  \mathit{g}%
_{yjj+r}-\frac{1}{3!}\mathit{g}_{yjj+r}^{3}+\frac{1}{5!}\mathit{g}_{yjj+r}%
^{5}-\cdots\right)  y_{j}y_{j+r}\right\}  ,
\end{align}
which can be further simplified as
\begin{equation}
\exp\left\{
%TCIMACRO{\tsum \limits_{1\leq j\leq r}}%
%BeginExpansion
{\textstyle\sum\limits_{1\leq j\leq r}}
%EndExpansion
\mathit{g}_{yjj+r}y_{j}y_{j+r}\right\}  =%
%TCIMACRO{\tprod \limits_{1\leq j\leq r}}%
%BeginExpansion
{\textstyle\prod\limits_{1\leq j\leq r}}
%EndExpansion
\left\{  \cos\left(  \mathit{g}_{yjj+r}\right)  I_{\mathcal{F}}+\sin\left(
\mathit{g}_{yjj+r}\right)  y_{j}y_{j+r}\right\}  .
\end{equation}
The expansion of the complex exponential gives
\begin{align}
\exp\left\{  ig\right\}   &  =I+ig+\frac{1}{2!}i^{2}g^{2}+i^{3}\frac{1}%
{3!}g^{3}+\cdots=\exp\left\{  i%
%TCIMACRO{\tsum \limits_{1\leq j\leq r}}%
%BeginExpansion
{\textstyle\sum\limits_{1\leq j\leq r}}
%EndExpansion
\mathit{g}_{yjj+r}y_{j}y_{j+r}\right\}  \\
&  =%
%TCIMACRO{\tprod \limits_{1\leq j\leq r}}%
%BeginExpansion
{\textstyle\prod\limits_{1\leq j\leq r}}
%EndExpansion
\left\{  \cos\left(  \mathit{g}_{yjj+r}\right)  I_{\mathcal{F}}+i\sin\left(
\mathit{g}_{yjj+r}\right)  y_{j}y_{j+r}\right\}  =%
%TCIMACRO{\tprod \limits_{1\leq j\leq r}}%
%BeginExpansion
{\textstyle\prod\limits_{1\leq j\leq r}}
%EndExpansion
\exp\left\{  i\mathit{g}_{yjj+r}y_{j}y_{j+r}\right\}  .
\end{align}


The values of $f$ are generated by
\begin{align}
&  \left.  \frac{\partial\Phi_{rD}\left(  g\right)  }{\partial\mathit{g}%
_{yjj+r}}\right\vert _{g=0}=\mathit{\operatorname*{tr}}\left(  y_{j}%
y_{j+r}D\right)  =f_{D}\left(  y_{j}y_{j+r}\right)  \\
&  \left.  =\right.  \left.  \frac{\partial\mathit{\operatorname*{tr}}\left(
%TCIMACRO{\tprod \limits_{1\leq j\leq r}}%
%BeginExpansion
{\textstyle\prod\limits_{1\leq j\leq r}}
%EndExpansion
\left\{  \cos\left(  \mathit{g}_{yjj+r}\right)  I_{\mathcal{F}}+\sin\left(
\mathit{g}_{yjj+r}\right)  y_{j}y_{j+r}\right\}  D\right)  }{\partial
\mathit{g}_{yjj+r}}\right\vert _{g=0}\\
&  \left.  \frac{\partial^{2}\Phi_{rD}\left(  g\right)  }{\partial
\mathit{g}_{yj_{1}j_{1}+r}\partial\mathit{g}_{yj_{2}j_{2}+r}}\right\vert
_{g=0}=\mathit{\operatorname*{tr}}\left(  y_{j_{1}}y_{j_{1}+r}y_{j_{2}%
}y_{j_{2}+r}D\right)  =f_{D}\left(  y_{j_{1}}y_{j_{1}+r}y_{j_{2}}y_{j_{2}%
+r}\right)  \\
&  \left.  =\right.  \left.  \frac{\partial^{2}\mathit{\operatorname*{tr}%
}\left(
%TCIMACRO{\tprod \limits_{1\leq j\leq r}}%
%BeginExpansion
{\textstyle\prod\limits_{1\leq j\leq r}}
%EndExpansion
\left\{  \cos\left(  \mathit{g}_{yjj+r}\right)  I_{\mathcal{F}}+\sin\left(
\mathit{g}_{yjj+r}\right)  y_{j}y_{j+r}\right\}  D\right)  }{\partial
\mathit{g}_{yj_{1}j_{1}+r}\partial\mathit{g}_{yj_{2}j_{2}+r}}\right\vert
_{g=0}\\
&  \vdots
\end{align}
and
\begin{align}
-i\left.  \frac{\partial\Phi_{iD}\left(  g\right)  }{\partial\mathit{g}%
_{yjj+r}}\right\vert _{g=0} &  =\mathit{\operatorname*{tr}}\left(
y_{j}y_{j+r}D\right)  =f_{D}\left(  y_{j}y_{j+r}\right)  =-i\left.
\frac{\partial%
%TCIMACRO{\tprod \limits_{1\leq j\leq r}}%
%BeginExpansion
{\textstyle\prod\limits_{1\leq j\leq r}}
%EndExpansion
\mathit{\operatorname*{tr}}\left(  \exp\left\{  i\mathit{g}_{yjj+r}%
y_{j}y_{j+r}\right\}  D\right)  }{\partial\mathit{g}_{yjj+r}}\right\vert
_{g=0}\\
-\left.  \frac{\partial^{2}\Phi_{rD}\left(  g\right)  }{\partial
\mathit{g}_{yj_{1}j_{1}+r}\partial\mathit{g}_{yj_{2}j_{2}+r}}\right\vert
_{g=0} &  =\mathit{\operatorname*{tr}}\left(  y_{j_{1}}y_{j_{1}+r}y_{j_{2}%
}y_{j_{2}+r}D\right)  =f_{D}\left(  y_{j_{1}}y_{j_{1}+r}y_{j_{2}}y_{j_{2}%
+r}\right)  =-\left.  \frac{\partial^{2}%
%TCIMACRO{\tprod \limits_{1\leq j\leq r}}%
%BeginExpansion
{\textstyle\prod\limits_{1\leq j\leq r}}
%EndExpansion
\mathit{\operatorname*{tr}}\left(  \exp\left\{  i\mathit{g}_{yjj+r}%
y_{j}y_{j+r}\right\}  D\right)  }{\partial\mathit{g}_{yj_{1}j_{1}+r}%
\partial\mathit{g}_{yj_{2}j_{2}+r}}\right\vert _{g=0}\\
&  \vdots
\end{align}


The characteristic function $\Psi_{rD}$ is \emph{positive semi definite} and
\emph{real positive semi definite if the geminals commute }i.e.
\begin{equation}
\sum_{1\leq j,k\leq m}\Psi_{rD}\left(  g_{j}-g_{k}\right)  c_{j}\bar{c}%
_{k}\geq0;1\leq m\leq\infty.
\end{equation}
This is a result of the factorization property%
\begin{equation}
\Phi_{rD}\left(  g_{j}-g_{k}\right)  =\mathit{\operatorname*{tr}}\left\{
\exp\left\{  \left(  g_{j}-g_{k}\right)  \right\}  D\right\}
=\mathit{\operatorname*{tr}}\left\{  \exp g_{j}\exp\left\{  -g_{k}\right\}
D\right\}  =\mathit{\operatorname*{tr}}\left\{  \exp g_{j}\left(  \exp\left\{
g_{k}\right\}  \right)  ^{\dagger}D\right\}  ,
\end{equation}
which leads to
\begin{equation}
\sum_{1\leq j,k\leq m}\Phi_{D}\left(  g_{j}-g_{k}\right)  c_{j}\bar{c}%
_{k}=\mathit{\operatorname*{tr}}\left\{  \left(  \sum_{1\leq j\leq m}c_{j}%
\exp\left\{  g_{j}\right\}  \right)  \left(  \sum_{1\leq k\leq m}c_{k}%
\exp\left\{  g_{k}\right\}  \right)  ^{\dagger}D\right\}
=\mathit{\operatorname*{tr}}\left\{  \Lambda\Lambda^{\dagger}D\right\}  ,
\end{equation}
where
\begin{equation}
\Lambda=\sum_{1\leq j\leq\infty}c_{j}\exp\left\{  g_{j}\right\}  .
\end{equation}
The characteristic function $\Psi_{iD}$ is \emph{real positive semi definite}
i.e.
\begin{equation}
\sum_{1\leq j,k\leq m}\Phi_{iD}\left(  g_{j}-g_{k}\right)  c_{j}c_{k}%
\geq0;1\leq m\leq\infty.
\end{equation}
This is a result of the factorization property%
\begin{equation}
\Phi_{iD}\left(  g_{j}-g_{k}\right)  =\mathit{\operatorname*{tr}}\left\{
\exp\left\{  \left(  g_{j}-g_{k}\right)  \right\}  D\right\}
=\mathit{\operatorname*{tr}}\left\{  \exp g_{j}\exp\left\{  -g_{k}\right\}
D\right\}  =\mathit{\operatorname*{tr}}\left\{  \exp g_{j}\left(  \exp\left\{
g_{k}\right\}  \right)  ^{t}D\right\}  ,
\end{equation}
which leads to
\begin{equation}
\sum_{1\leq j,k\leq m}\Phi_{D}\left(  g_{j}-g_{k}\right)  c_{j}c_{k}%
=\mathit{\operatorname*{tr}}\left\{  \left(  \sum_{1\leq j\leq m}c_{j}%
\exp\left\{  g_{j}\right\}  \right)  \left(  \sum_{1\leq k\leq m}c_{k}%
\exp\left\{  g_{k}\right\}  \right)  ^{t}D\right\}
=\mathit{\operatorname*{tr}}\left\{  \Lambda\Lambda^{t}D\right\}  ,
\end{equation}
where
\begin{equation}
\Lambda=\sum_{1\leq j\leq\infty}c_{j}\exp\left\{  g_{j}\right\}  .
\end{equation}


$\lceil$One could also consider the\emph{ exterior exponential }%
$\operatorname{Exp}\left\{  A\right\}  $ which is defined as
\begin{equation}
\operatorname{Exp}\left\{  A\right\}  =I+A+\frac{1}{2!}A\wedge A+\frac{1}%
{3!}A\wedge A\wedge A+\cdots
\end{equation}
to construct another characteristic function $\Psi_{D}\left(  g\right)  $,
which is $\emph{gauge}$ $\emph{invariant}$, can be obtained by using the
exterior exponential of a geminal
\begin{equation}
\operatorname{Exp}\left\{  ig\right\}  =I_{\mathcal{F}}+g+\frac{1}{2!}%
i^{2}g\wedge g+\frac{1}{3!}i^{3}g\wedge g\wedge g+\cdots,
\end{equation}
however $\operatorname{Exp}$ does not have the desired factorization
properties for positive definiteness.$\rfloor$

\paragraph{Creator Form of Geminal Operators}

A restricted form of a geminal operator is given by
\begin{equation}
g=%
%TCIMACRO{\tsum \limits_{1\leq j_{1}<j_{2}\leq2r}}%
%BeginExpansion
{\textstyle\sum\limits_{1\leq j_{1}<j_{2}\leq2r}}
%EndExpansion
\mathit{g}_{j_{1}j_{2}}a_{j_{1}}^{\dagger}a_{j_{2}}^{\dagger}%
\end{equation}
here
\begin{equation}
\left[  a_{j_{1}}^{\dagger}a_{j_{2}}^{\dagger},a_{k_{1}}^{\dagger}a_{k_{2}%
}^{\dagger}\right]  =0
\end{equation}
and
\begin{equation}
g=-g^{t}.
\end{equation}
Thus
\begin{equation}
e^{g}e^{g^{t}}=I_{\mathcal{F}}%
\end{equation}


\subsubsection{Cumulants}%

\begin{equation}
\Omega\left(  g\right)  =\sum_{0\leq m\leq2r}\frac{1}{m!}%
%TCIMACRO{\tsum \limits_{1\leq\left(  j_{1}<j_{2}\right)  ,\vdots\cdots,\left(
%j_{m-1}<j_{m}\right)  \leq2r}}%
%BeginExpansion
{\textstyle\sum\limits_{1\leq\left(  j_{1}<j_{2}\right)  ,\vdots\cdots,\left(
j_{m-1}<j_{m}\right)  \leq2r}}
%EndExpansion
\mathit{g}_{j_{1}j_{2}}\cdots\mathit{g}_{j_{m-1}j_{m}}\left.  \frac
{\partial^{m}\Omega\left(  g\right)  }{\partial\mathit{g}_{j_{1}j_{2}}%
\cdots\partial\mathit{g}_{j_{m-1}j_{m}}}\right\vert _{g=0}%
\end{equation}
%

\begin{align}
\Omega\left(  g\right)   &  =\operatorname{Log}\left\{  \Phi_{D}\left(
g\right)  \right\}  =\sum_{0\leq m\leq2r}%
%TCIMACRO{\tsum \limits_{1\leq\left(  j_{1}<j_{2}\right)  ,\vdots\cdots,\left(
%j_{m-1}<j_{m}\right)  \leq2r}}%
%BeginExpansion
{\textstyle\sum\limits_{1\leq\left(  j_{1}<j_{2}\right)  ,\vdots\cdots,\left(
j_{m-1}<j_{m}\right)  \leq2r}}
%EndExpansion
P_{j_{1}\cdots j_{m}}\left(  g\right)  \mathit{g}_{j_{1}j_{2}}\cdots
\mathit{g}_{j_{m-1}j_{m}}\\
&  =\sum_{0\leq m\leq2r}\frac{1}{m!}%
%TCIMACRO{\tsum \limits_{1\leq\left(  j_{1}<j_{2}\right)  ,\vdots\cdots,\left(
%j_{m-1}<j_{m}\right)  \leq2r}}%
%BeginExpansion
{\textstyle\sum\limits_{1\leq\left(  j_{1}<j_{2}\right)  ,\vdots\cdots,\left(
j_{m-1}<j_{m}\right)  \leq2r}}
%EndExpansion
\mathit{g}_{j_{1}j_{2}}\cdots\mathit{g}_{j_{m-1}j_{m}}\left.  \frac
{\partial^{m}\operatorname{Log}\left\{  \Phi_{D}\left(  g\right)  \right\}
}{\partial\mathit{g}_{j_{1}j_{2}}\cdots\partial\mathit{g}_{j_{m-1}j_{m}}%
}\right\vert _{g=0}\\
\Phi_{D}\left(  g\right)   &  =\exp\left\{  \Omega\left(  g\right)  \right\}
=\exp\left\{  \sum_{0\leq m\leq2r}\frac{1}{m!}%
%TCIMACRO{\tsum \limits_{1\leq\left(  j_{1}<j_{2}\right)  ,\vdots\cdots,\left(
%j_{m-1}<j_{m}\right)  \leq2r}}%
%BeginExpansion
{\textstyle\sum\limits_{1\leq\left(  j_{1}<j_{2}\right)  ,\vdots\cdots,\left(
j_{m-1}<j_{m}\right)  \leq2r}}
%EndExpansion
\mathit{g}_{j_{1}j_{2}}\cdots\mathit{g}_{j_{m-1}j_{m-2}}\left.  \frac
{\partial^{m}\operatorname{Log}\left\{  \Phi_{D}\left(  g\right)  \right\}
}{\partial\mathit{g}_{j_{1}j_{2}}\cdots\partial\mathit{g}_{j_{m-1}j_{m}}%
}\right\vert _{g=0}\right\}  .
\end{align}%
\begin{equation}
\operatorname{Log}\left\{  \Phi_{D_{r}}\left(  g\right)  \right\}
=\operatorname{Log}\left\{  \mathit{\operatorname*{tr}}\left(  \exp\left\{
g\right\}  D\right)  \right\}  =\operatorname{Log}\left\{
\mathit{\operatorname*{tr}}\left(
%TCIMACRO{\tprod \limits_{1\leq j_{1}<j_{2}\leq2r}}%
%BeginExpansion
{\textstyle\prod\limits_{1\leq j_{1}<j_{2}\leq2r}}
%EndExpansion
\left\{  \cos\left(  \mathit{g}_{j_{1}j_{2}}\right)  I_{\mathcal{F}}%
+\sin\left(  \mathit{g}_{j_{1}j_{2}}\right)  x_{j_{1}}x_{j_{2}}\right\}
D\right)  \right\}
\end{equation}%
\begin{equation}
\operatorname{Log}\left\{  \Phi_{D_{i}}\left(  g\right)  \right\}
=\operatorname{Log}\left\{  \mathit{\operatorname*{tr}}\left(  \exp\left\{
g\right\}  D\right)  \right\}  =\operatorname{Log}\left\{
\mathit{\operatorname*{tr}}\left(
%TCIMACRO{\tprod \limits_{1\leq j_{1}<j_{2}\leq2r}}%
%BeginExpansion
{\textstyle\prod\limits_{1\leq j_{1}<j_{2}\leq2r}}
%EndExpansion
\exp\left\{  i\mathit{g}_{j_{1}j_{2}}x_{j_{1}}x_{j_{2}}\right\}  D\right)
\right\}
\end{equation}



\end{document}